%!TEX root=../mythesis.tex
% Chapter 3

\chapter{Methodology}
\chaptermark{Methodology}
\label{ch:methodology}

\section{Hybrid Biomass Power-to-Methanol System Model}

\begin{itemize}
\item Energy system definition.
\item Technical parameters definition.
\item Economical parameters definition.
\end{itemize}

\section{Two-steps Outerloop \& Innerloop Strategy}

The methodology adopts a two-step strategy in which layout/sizing and operational dispatch decisions are decomposed into outerloop and innerloop problems to improve tractability and interpretability.

\section{Level 2 Optimization Architecture (MILP)}

\begin{itemize}
\item Component-level non-linear modeling logic.
\item Operation logic.
\item Modeling tool: \texttt{oemof.solph}.
\item Solver: CBC.
\item Objective function definition.
\end{itemize}

\section{Level 3 Optimization Architecture (RL)}

\begin{itemize}
\item Component-level linear modeling logic.
\item Chosen RL algorithm.
\item Operation logic (MDP).
\item Integration of Bayesian Optimization (BO) and RL.
\item Objective function definition (reward function).
\end{itemize}

\section{Data}

This section is reserved for dataset sources, preprocessing, assumptions, and scenario-specific inputs.

\section{Comparison Stages \& Scenario Definition}

This comparison is conducted in two steps:

\begin{itemize}
\item \textbf{Outerloop comparison}
\begin{itemize}
\item Innerloop operational decision is fixed as an \texttt{oemof} problem.
\item For outerloop (layout selection fed into innerloop), different tools are compared: \texttt{oemof} investment block (IB), Bayesian Optimization (BO), and Design of Experiment (DOE).
\item Objective: identify which approach gives the best layout under both linear and non-linear economies of scale.
\item Components to optimize: electricity supply and Power-to-X blocks.
\end{itemize}

\item \textbf{Innerloop comparison}
\begin{itemize}
\item Compare RL dispatch in a physics-based plant model against \texttt{oemof} dispatch in a linear model.
\item Stage 1 (fixed outerloop, linear model with RL): evaluate how close RL can optimize a linear model compared to Level 2 MILP.
\item Stage 2 (fixed outerloop, non-linear model with RL): evaluate behavior when optimized linear layouts are injected into a physics-based model, including feasibility and performance changes.
\item Stage 3 (joint outerloop + innerloop): compare BO+RL non-linear pathway vs BO+\texttt{oemof} linear pathway.
\end{itemize}
\end{itemize}

Case studies for each step (especially innerloop stage comparisons):

\begin{itemize}
\item \textbf{Weather variability}
\begin{itemize}
\item Stable solar and wind input.
\item Extreme weather volatility.
\end{itemize}
\item \textbf{Market conditions}
\begin{itemize}
\item Cheap batteries.
\item Cheap hydrogen system (e.g., -50\%, -20\%, +20\%, +50\%).
\end{itemize}
\item \textbf{System design and flexibility}
\begin{itemize}
\item Heavy fixed methanol reactor load.
\item Heavy fixed electrolyzer load.
\end{itemize}
\end{itemize}

For each step, stage, and case study, predetermined KPIs are examined:

\begin{itemize}
\item \textbf{Techno-economic level}
\begin{itemize}
\item LCOM.
\item CAPEX and OPEX.
\item Carbon intensity.
\item Grid dependency ratio.
\item Self-consumption ratio.
\item Autonomy days.
\end{itemize}
\item \textbf{Framework level}
\begin{itemize}
\item Simulation time.
\item BO distance.
\end{itemize}
\item \textbf{Scenario descriptors}
\begin{itemize}
\item Year.
\item CAPEX and OPEX per component.
\item Conversion factor per component.
\end{itemize}
\item \textbf{Fixed components}
\begin{itemize}
\item MeOH reactor.
\item Distillery.
\item MeOH storage.
\item Anaerobic digester.
\item ATR.
\item O$_2$ storage (full ATR-load sizing for 3 days).
\item Steam generator.
\item Auxiliary units: RWGS, fuel boiler, AD gas separator, compressor, and three heat exchangers.
\end{itemize}
\item \textbf{Variable components (IB vs BO vs DOE)}
\begin{itemize}
\item CH$_4$ generator or gas turbine.
\item Wind turbine.
\item PV.
\item BESS.
\item Electrolyzer.
\item Hydrogen storage.
\item Biogas storage.
\end{itemize}
\end{itemize}
